\documentclass[12pt,a4paper]{article}

\usepackage{amsmath,amssymb,amsfonts}
\usepackage{geometry}
\usepackage[colorlinks=true,linkcolor=red!70!black,citecolor=green!50!black,urlcolor=blue!80!black]{hyperref}
\usepackage{booktabs}
\usepackage{physics}
\usepackage{xcolor}
\usepackage{graphicx}
\usepackage[numbers]{natbib}

\geometry{margin=2.5cm}


% ─── Macros SSEE ────────────────────────────────────────────────────────────
\newcommand{\phiG}{\varphi}
\newcommand{\KAL}{\mathrm{KAL}_0}
\newcommand{\KALeff}{\mathrm{KAL}_{\rm eff}}
\newcommand{\MIRA}{\mathcal{M}}
\newcommand{\AURA}{\mathcal{A}}
\newcommand{\Mpl}{M_{\mathrm{Pl}}}
\newcommand{\Omde}{\Omega_{\mathrm{DE}}}
\newcommand{\Omm}{\Omega_{m}}
\newcommand{\half}{\tfrac{1}{2}}
\newcommand{\Hzero}{H_0}
\newcommand{\bc}{\beta_c}
\newcommand{\fsc}{f_{\mathrm{screen}}}
\newcommand{\sK}{s_K}
\renewcommand{\Omega}{\varOmega}
\newcommand{\Mcuv}{M^4_{\rm UV}}
\newcommand{\rhoc}{\rho_{\rm crit}}
% ─────────────────────────────────────────────────────────────────────────────

\title{\textbf{UV Extension of SSEE Dark Energy:\\
       $M^4 = 45\alpha^2\rho_{\rm crit} = 5\phiG^8\rho_{\rm crit}$\\
       and a Conditional Self-Consistency Check on the Hubble Tension}}

\author{Mike Edison Almeida Vallejo \\
\small \textit{Independent Researcher} \\
\small \href{mailto:mike.almeida1721@gmail.com}{mike.almeida1721@gmail.com} \\
\small ORCID: \href{https://orcid.org/0009-0008-2195-7836}{0009-0008-2195-7836}
}

\date{May 2026 --- Paper 10 of the SSEE series}

\sloppy
% ACTA-PROCEDENCIA {"commit": "5651fc97a90dc17d52be0317b46142db9ac86fcb", "reproducible_desde_commit": true, "cambios_sin_commitear": [], "script": "src/valores_tex.py", "script_sha256": "cb5ceb2369bdef5e586d471d57d44a11a9fed354be540a96ce34da85e582efa4", "nucleo_sha256": "5cec8af72d88f154116c5a532e42a9b6ff7a5b4c3d86bcc6921eabb246698993", "entradas": {"manuscript/valores.yaml": "5fa01e91dd68164243d664a73264bb6c87431d2ec0bd69db87c62a983a55992d", "results/logs/kids_publicados.json": "6c25c98059ce91038215358b4e7ac9cfb40198f410354937cee4ea00b46654ff", "results/logs/p10_uv_fscreen.json": "8fb68b295c848e0292e5701e7de795ba7c7ba25e3292d776260bb3dea8970c95", "results/logs/s8_desde_b1.json": "6d225f4403fbfbe6934d7edca26aaefbbe19e05ef6c7b7eda842a259e798b1c1"}, "argv": ["src/valores_tex.py"], "fecha": "2026-09-30T19:35:24", "host": "mike-System-Product-Name", "python": "3.12.3", "versiones": {}, "nucleos": {"OMP_NUM_THREADS": null, "MKL_NUM_THREADS": null, "OPENBLAS_NUM_THREADS": null, "OMPI_COMM_WORLD_SIZE": null}}
% GENERADO por src/valores_tex.py desde manuscript/valores.yaml — NO EDITAR A MANO.
% Cada numero sale de un log de la cadena de procedencia (dvc.yaml + acta).
\providecommand{\val}[1]{\csname ssee@val@#1\endcsname}
\expandafter\def\csname ssee@val@H0_glob\endcsname{67.962142}  % results/logs/p10_uv_fscreen.json#H0_glob
\expandafter\def\csname ssee@val@PTE_k1000_lcdm\endcsname{0.010}  % results/logs/kids_publicados.json#PTE.lcdm.PTE
\expandafter\def\csname ssee@val@PTE_k1000_ssee\endcsname{0.012}  % results/logs/kids_publicados.json#PTE.ssee.PTE
\expandafter\def\csname ssee@val@S8_b1_lcdm\endcsname{0.8229}  % results/logs/s8_desde_b1.json#filas.LCDM.S8
\expandafter\def\csname ssee@val@S8_b1_lcdm_sigma\endcsname{0.0059}  % results/logs/s8_desde_b1.json#filas.LCDM.S8_sigma
\expandafter\def\csname ssee@val@S8_b1_ssee\endcsname{0.8261}  % results/logs/s8_desde_b1.json#filas.SSEE.S8
\expandafter\def\csname ssee@val@S8_b1_ssee_sigma\endcsname{0.0054}  % results/logs/s8_desde_b1.json#filas.SSEE.S8_sigma
\expandafter\def\csname ssee@val@S8_legacy_dato\endcsname{0.8265}  % results/logs/kids_publicados.json#kids_legacy.S8
\expandafter\def\csname ssee@val@S8_legacy_dato_sigma\endcsname{0.0176}  % results/logs/kids_publicados.json#kids_legacy.S8_sigma
\expandafter\def\csname ssee@val@chi2_k1000_lcdm\endcsname{262.75}  % results/logs/kids_publicados.json#PTE.lcdm.chi2
\expandafter\def\csname ssee@val@chi2_k1000_publicado\endcsname{260.31}  % results/logs/kids_publicados.json#kids1000_maxpost.chi2
\expandafter\def\csname ssee@val@chi2_k1000_ssee\endcsname{265.44}  % results/logs/kids_publicados.json#PTE.ssee.chi2
\expandafter\def\csname ssee@val@dAIC_k1000\endcsname{-5.31}  % results/logs/kids_publicados.json#seleccion_kids1000.delta_AIC
\expandafter\def\csname ssee@val@dBIC_k1000\endcsname{-19.0}  % results/logs/kids_publicados.json#seleccion_kids1000.delta_BIC
\expandafter\def\csname ssee@val@dchi2_k1000\endcsname{2.69}  % results/logs/kids_publicados.json#seleccion_kids1000.delta_chi2
\expandafter\def\csname ssee@val@fscreen_IR\endcsname{0.067253}  % results/logs/p10_uv_fscreen.json#f_screen_IR
\expandafter\def\csname ssee@val@fscreen_UV_corr\endcsname{0.002268}  % results/logs/p10_uv_fscreen.json#f_screen_UV_corr
\expandafter\def\csname ssee@val@fscreen_full\endcsname{0.069522}  % results/logs/p10_uv_fscreen.json#f_screen_full
\expandafter\def\csname ssee@val@logA_b1_lcdm\endcsname{3.045}  % results/logs/s8_desde_b1.json#filas.LCDM.logA
\expandafter\def\csname ssee@val@logA_b1_ssee\endcsname{3.042}  % results/logs/s8_desde_b1.json#filas.SSEE.logA
\expandafter\def\csname ssee@val@logA_b1_ssee_sigma\endcsname{0.013}  % results/logs/s8_desde_b1.json#filas.SSEE.logA_sigma
\expandafter\def\csname ssee@val@p_dchi2_k1000\endcsname{0.61}  % results/logs/kids_publicados.json#seleccion_kids1000.p_delta_chi2
\expandafter\def\csname ssee@val@sK_IR\endcsname{0.403302}  % results/logs/p10_uv_fscreen.json#s_K_IR
\expandafter\def\csname ssee@val@sK_UV_corr\endcsname{0.013603}  % results/logs/p10_uv_fscreen.json#s_K_UV_corr
\expandafter\def\csname ssee@val@sK_full\endcsname{0.416905}  % results/logs/p10_uv_fscreen.json#s_K_full
\expandafter\def\csname ssee@val@tension_b1_lcdm\endcsname{2.59}  % results/logs/s8_desde_b1.json#filas.LCDM.tension_kids
\expandafter\def\csname ssee@val@tension_b1_ssee\endcsname{2.73}  % results/logs/s8_desde_b1.json#filas.SSEE.tension_kids
\expandafter\def\csname ssee@val@z_k1000_lcdm\endcsname{2.5}  % results/logs/kids_publicados.json#PTE.lcdm.z
\expandafter\def\csname ssee@val@z_k1000_ssee\endcsname{2.4}  % results/logs/kids_publicados.json#PTE.ssee.z

\begin{document}
\maketitle

\begin{abstract}
Within the Structural Self-Energy Expansion (SSEE), the infrared (IR) kinetic
Lagrangian $K(X) = X/\KAL$ produces a local Hubble correction
$H_0^{\rm glob} = H_0^{\rm local}(1-f_{\rm scr})$ with $f_{\rm scr}\approx 0.06725$
(Paper~9). Run on the only directly measured rate, $H_0^{\rm SH0ES}=73.04\pm1.04$,
this gives $H_0^{\rm glob,IR}=68.13$ km/s/Mpc, $0.17\sigma$ from the dimensionless
algebraic number $3(\phiG+\pi)^2=67.96214$ (Postulate~D, Paper~1); the
earlier MIRA-mapped anchor $67.037$ is superseded
by the 2026 reframe.
We extend the EFT to include the leading higher-derivative operator
$K(X) = X/\KAL + X^2/M^4$ and identify the unique algebraic value of the
UV cutoff $M$ such that the theory remains connected to the inflationary
$\alpha$-attractor of Paper~1.
The $\alpha$-attractor parameter $\alpha = \phiG^4/3$, which fixes the
inflationary tensor-to-scalar ratio $r = 12\alpha/N^2 \approx 0.00813$
(a pre-committed LiteBIRD target), also determines
\begin{equation*}
  \boxed{M^4 \;=\; 45\,\alpha^2\,\rhoc \;=\; 5\,\phiG^8\,\rhoc
         \;\approx\; 234.9\,\rhoc}\,,
\end{equation*}
where the identity $45\alpha^2 = 5\phiG^8$ follows exactly from the definition
$\alpha = \phiG^4/3$ (not a numerical coincidence, but an algebraic consequence).
With this cutoff, the full screening rate becomes
$\sK^{\rm full} = \val{S_K_FULL_d5}$ ($+3.37\%$ over the IR value). Run on the SH0ES
datum $73.04\pm1.04$ km/s/Mpc \citep{Riess2022}, the cascade returns
$H_0^{\rm glob,UV} = 67.96214$ km/s/Mpc, which differs from the pure number
$3(\phiG+\pi)^2 = 67.96214$ by $4.2\times10^{-6}$ --- the headline UV
result of this paper, refining the IR residual $0.166$ ($0.17\sigma$).
\textbf{The UV value is conditional on Postulate~C.1 (UV-IR $\varphi/\pi$
separability; §\ref{sec:theorem})}, which fixes $\sK^{\rm full}$ via the
cutoff $M$; the IR value is independent of it.
The same $\alpha$ that governs inflation at $z \sim 10^{60}$ thus fixes the
dark-energy UV scale today, providing a falsifiable cross-epoch consistency test:
if LiteBIRD measures $r \neq 0.00813$, both $\alpha$ and $M$ change, and the
predicted $H_0^{\rm glob}$ changes accordingly.  The IR result
$H_0^{\rm glob,IR} = 68.13$ km/s/Mpc ($0.17\sigma$; Paper~9) is
independent of this postulate and stands on its own.
\end{abstract}

\clearpage
{\small\tableofcontents}
\clearpage

% ═══════════════════════════════════════════════════════════════════════════════
\section{Introduction}
% ═══════════════════════════════════════════════════════════════════════════════

The discovery of cosmic acceleration \citep{Riess1998,Perlmutter1999} motivated
scalar-field dark energy models ranging from quintessence \citep{Ratra1988,Caldwell2002}
to k-essence \citep{ArmendarizPicon2001} and general Horndeski theories
\citep{Horndeski1974,Copeland2006,Weinberg2013review}.
The Hubble tension---the $\sim4$--$5\sigma$ discrepancy between CMB-inferred
$H_0 = 67.36 \pm 0.54$ km/s/Mpc \citep{Planck2018} and local distance-ladder
measurements \citep{Verde2019,Kamionkowski2023}---has
resisted resolution through conventional extensions of $\Lambda$CDM.
Paper~9 of this series~\citep{SSEE9} showed that within SSEE, the dark-energy
pressure-dilution rate $\sK = -3w_0(1+w_0) = \val{S_K_d6}$ \citep{SSEE7,SSEE8} and the lensing factor
$\MIRA = \AURA/2$ combine to produce an algebraic screening fraction
\begin{equation}
  \fsc \;=\; \frac{\sK}{3\,\MIRA}
       \;=\; \frac{\pi - \phiG}{(\phiG+\pi)^2}
       \;=\; 0.06725\,,
\end{equation}
which converts the measured local rate into a global one,
\begin{equation}
  H_0^{\rm glob,IR} \;=\; H_0^{\rm SH0ES}\,(1 - \fsc) \;=\; 68.13 \text{ km/s/Mpc}\,,
\end{equation}
to be compared with $3(\phiG+\pi)^2 = 67.96214$ --- a residual of $0.166$,
i.e.\ $0.17\sigma$ on the propagated uncertainty $1.04\times(1-\fsc)=0.970$.
The comparison target is a pure number and carries no units; it is therefore
never the input of the cascade, only its benchmark. The earlier
register used the MIRA-mapped CMB anchor $H_0^{\rm MIRA} = 67.037$ ($\to 71.87$,
$1.12\sigma$); this is now superseded. Under the $\omega_m$-direct reframe the
CMB matter density is the standard $\Omega_{m,\mathrm{CMB}}=\omega_m/h^2=\val{OMEGA_M_TOTAL_d6}$
(no matter factor; OP-8 dissolved), built from algebraic $\omega_b,\omega_c$, and
the single number $67.96214$ matches both the cascade output and the CMB-preferred value
(Paper~3).

That IR result uses the kinetic Lagrangian $K(X) = X/\KAL$ (the leading term),
and the UV cutoff $M \to \infty$ is implicit.
This paper asks: \emph{what is $M$ physically?}

The leading higher-derivative correction is $K(X) = X/\KAL + X^2/M^4$.
We show that the unique algebraically consistent value of $M$ is
\begin{equation}
  M^4 \;=\; 45\,\alpha^2\,\rhoc \;=\; 5\,\phiG^8\,\rhoc\,,
  \label{eq:M4main}
\end{equation}
where $\alpha = \phiG^4/3$ is the inflationary $\alpha$-attractor parameter of
Paper~1~\citep{SSEE1}, and the equality $45\alpha^2 = 5\phiG^8$ is an exact algebraic identity
following from the definition of $\alpha$.
With this value, the global rate returned from the SH0ES datum
$73.04\pm1.04$ km/s/Mpc \citep{Riess2022} is
$H_0^{\rm glob,UV} = \val{H0_GLOBAL_d5}$ km/s/Mpc, coincident with $3(\phiG+\pi)^2$ to
$4.2\times10^{-6}$ and refining the IR residual $0.166$ ($0.17\sigma$). The UV value
is conditional on Postulate~C.1 (see §\ref{sec:theorem}), which fixes
$\sK^{\rm full}$ via the cutoff $M$; the IR value is independent of it.

The physical picture is intriguing: the \emph{same curvature parameter $\alpha$
that shaped the inflationary plateau at $z \sim 10^{60}$} provides the speed
limit $M$ for the dark-energy kinetic term today.
Section~\ref{sec:cross} makes this a falsifiable cross-epoch prediction.


% ═══════════════════════════════════════════════════════════════════════════════
\section{The $\alpha$-Attractor Parameter and the UV Identity}
\label{sec:alpha}
% ═══════════════════════════════════════════════════════════════════════════════

\subsection{The inflaton curvature $\alpha$}

Paper~1 of this series derives an inflationary $\alpha$-attractor potential
with curvature parameter
\begin{equation}
  \alpha \;=\; \frac{\phiG^4}{3}
         \;\approx\; \val{ALPHA_ATT_d4}\,,
  \label{eq:alpha}
\end{equation}
where $\phiG = (1+\sqrt{5})/2$ is the golden ratio.
In the Starobinsky-like limit \citep{Starobinsky1980,Kallosh2013,Galante2015}, this predicts a tensor-to-scalar ratio
\begin{equation}
  r \;=\; \frac{12\alpha}{N^2}
    \;\approx\; \frac{12 \times \val{ALPHA_ATT_d4}}{58.07^2}
    \;\approx\; 0.00813
  \label{eq:r}
\end{equation}
for $N = 2\varphi^7 \approx 58.07$ e-folds, the algebraic e-fold count of
SSEE inflation (Paper~1), consistent with the standard window \citep{LiddleLeach2003}.
This prediction lies below the current upper bound $r < 0.036$
\citep{Akrami2020} and is within the anticipated sensitivity of LiteBIRD
\citep{LiteBIRD2020} (target $\sigma_r \sim 0.002$, expected results 2032)
and a future CMB-S4 survey \citep{CMBS4_2016}, making
Eq.~\eqref{eq:r} a pre-committed falsification criterion.

\subsection{The algebraic identity}

From the definition $\alpha = \phiG^4/3$:
\begin{align}
  45\,\alpha^2
    &= 45\,\left(\frac{\phiG^4}{3}\right)^2
     = 45\,\frac{\phiG^8}{9}
     = 5\,\phiG^8\,.
  \label{eq:identity}
\end{align}
This is an exact algebraic consequence of the definition $\alpha = \phiG^4/3$,
not a numerical coincidence: floating-point evaluation confirms $|45\alpha^2 - 5\phiG^8| = 0$
to double precision, as expected for a closed-form algebraic identity.

The factor 5 in $5\phiG^8$ is the same integer that generates the golden ratio,
$\phiG = (1+\sqrt{5})/2$, so $5 = (2\phiG - 1)^2$.
Equivalently,
\begin{equation}
  M \;=\; \phiG^2 \times 5^{1/4} \times \rhoc^{1/4}
    \;\approx\; 9.68 \text{ meV}\,,
  \label{eq:Mmev}
\end{equation}
where $\rhoc^{1/4} \approx 2.47$ meV is the SSEE critical-density scale today,
evaluated at the model's physical Hubble rate $H_0 = 67.96214$~km/s/Mpc
(the CMB-anchored value of Paper~III).\footnote{Convention note: in the
dimensionless formulation of §\ref{sec:EFT} we work in units $\rhoc = 1$, where
$M^4 = 5\phiG^8$ is a pure number ($\approx \val{_M4_UV_d2}$) and all downstream
observables ($s_K$, $f_{\rm screen}$, $\varepsilon$) are dimensionless
ratios. The conversion to physical units in Eq.~\eqref{eq:Mmev} uses the SSEE
$\rho_{\rm crit}$ consistently with the identity $M^4 = 5\phiG^8\,\rhoc$ of
Table~\ref{tab:identity}; we do \emph{not} borrow the external $\Lambda$CDM
scale $\rho_\Lambda^{1/4}\approx 2.25$~meV (which would give the inconsistent
$M\approx 8.81$~meV). Crucially, the meV value of $M$ is a display scale only:
the Hubble cascade of §\ref{sec:hubble} runs on the dimensionless ratio
$M^4/\rhoc = 5\phiG^8 = \val{_M4_UV_d6}$, so the UV Hubble prediction
($H_0^{\rm glob,UV}=\val{H0_GLOBAL_d5}$ from the SH0ES datum) is
unchanged under any choice of reference density.}

Table~\ref{tab:identity} summarises the equivalent forms of the identity.

\begin{table}[ht]
\centering
\caption{Equivalent algebraic forms of the UV cutoff $M^4$.}
\label{tab:identity}
\begin{tabular}{lll}
\toprule
Form & Expression & Numerical value \\
\midrule
$\alpha$-form & $M^4 = 45\,\alpha^2\,\rhoc$ & $\val{_M4_UV_d2}\,\rhoc$ \\
$\phiG$-form  & $M^4 = 5\,\phiG^8\,\rhoc$  & $\val{_M4_UV_d2}\,\rhoc$ \\
Root form     & $M = \phiG^2\,5^{1/4}\,\rhoc^{1/4}$ & $9.68$ meV \\
$(2\phiG{-}1)^2$ form & $M^4 = (2\phiG{-}1)^2 \phiG^8\,\rhoc$ & $\val{_M4_UV_d2}\,\rhoc$ \\
\bottomrule
\end{tabular}
\end{table}


% ═══════════════════════════════════════════════════════════════════════════════
\section{Full K(X) EFT and the Corrected $\sK$}
\label{sec:EFT}
% ═══════════════════════════════════════════════════════════════════════════════

\subsection{The effective description of the screening sector}

The screening rate itself is fixed by the equation of state alone
(Eq.~\eqref{eq:sKdef}) and is therefore independent of any Lagrangian. Its
\emph{UV correction} is not: computing it requires knowing what cuts the
theory off. We describe the sector that carries the screening by the
k-essence functional \citep{ArmendarizPicon2001}
\begin{equation}
  K(X) \;=\; \frac{X}{\KAL} + \frac{X^2}{M^4}\,,
  \label{eq:KX}
\end{equation}
where $X = -\tfrac{1}{2}g^{\mu\nu}\partial_\mu\phi\,\partial_\nu\phi$,
$\KAL = (\phiG+3\pi)/2 \approx 5.521$ is the SSEE structural retention
(Papers 1--7), and $M^4 = 45\alpha^2\rhoc$ is the UV cutoff introduced in
Eq.~\eqref{eq:M4main}.
The IR limit $M \to \infty$ recovers $K(X) = X/\KAL$, which reproduces
$\sK^{\rm IR} = \val{S_K_d6}$ exactly and is the description used in
Papers~1--9.

\paragraph{Units.}
Throughout this section we work in units $\rhoc = 1$. In those units $X$,
$K$ and $M^4$ all carry dimensions of energy density, so every UV correction
below is the dimensionless ratio $X^2/(M^4\rhoc)$ and \emph{not} $X^2/M^4$.
The distinction matters: with $\rhoc$ restored explicitly, the screening
fraction of Eq.~\eqref{eq:fUV} is invariant under a rescaling of $\rhoc$ over
$58$ orders of magnitude, returning $\fsc = \val{F_SCREEN_d12}$ at
$\rhoc = 1$, $10^{-47}$ and $3.7\times10^{11}$ alike. The cutoff therefore
enters only as a ratio, and no factor of $H_0$ reaches $\fsc$ through $M$.

\paragraph{This is not the dark-energy action of Paper~7.}
Paper~7 describes the dark-energy field as a ghost condensate,
$K = c_1 X + c_2 X^2$ with no potential and $c_2 X/c_1 = \val{u_condensado_d6}$; in
Eq.~\eqref{eq:KX} the same ratio is $+0.008577$. The two functionals do
different jobs: Paper~7's action reproduces $w_0$ and $c_s^2$ from the
Lagrangian, whereas Eq.~\eqref{eq:KX} transports the cutoff $M$ into the
screening rate. They are not interchangeable --- evaluating
Eq.~\eqref{eq:sKfullform} below with Paper~7's coefficients returns
$18.949$ instead of $\val{S_K_FULL_d5}$, a factor of $45.5$. Whether one theory
should furnish both is exactly what Postulate~C.1 leaves open
(§\ref{sec:firstprinciples}); nothing in the cascade of
§\ref{sec:hubble} depends on the answer, because $\sK^{\rm IR}$ comes
from $w_0$ and the correction is measured against the cutoff, not
against $c_2$.

\subsection{Self-consistent background kinetic energy}

The background equation of motion for $K(X) = X/\KAL + X^2/M^4$ is
\begin{equation}
  2X K_X \;=\; \rho_\phi(1+w_0)\,,
  \quad
  K_X \;=\; \frac{1}{\KAL} + \frac{2X}{M^4}\,,
\end{equation}
giving a quadratic equation for the self-consistent $X_{\rm bg}$:
\begin{equation}
  \frac{4\,X_{\rm bg}^2}{M^4} + \frac{2\,X_{\rm bg}}{\KAL}
  - \rho_\phi(1+w_0) \;=\; 0\,.
  \label{eq:quadratic}
\end{equation}
The physical (positive) root is
\begin{equation}
  X_{\rm bg}^{\rm UV}
  \;=\; \frac{-1/\KAL + \sqrt{1/\KAL^2 + 4\,\rho_\phi(1+w_0)/M^4}}
             {4/M^4}\,.
\end{equation}
With $M^4 = 45\alpha^2\rhoc = \val{_M4_UV_d2}\rhoc$ and the SSEE background
$\rho_\phi(1+w_0) = \sK^{\rm IR}/3 = \val{sK_tercio_d5}$ (in units $\rhoc = 1$),
we find
\begin{equation}
  X_{\rm bg}^{\rm UV} \;=\; \val{_X_UV_d5}\,\rhoc\,,
  \quad
  \frac{X_{\rm bg}^{\rm UV}}{X_{\rm bg}^{\rm IR}} \;=\; 0.9831\,,
\end{equation}
a 1.7\% reduction from the IR value $X_{\rm bg}^{\rm IR} = \val{X_bg_IR_d5}\,\rhoc$.
The UV correction is perturbative: $\varepsilon \equiv X_{\rm bg}^2/M^4
\approx 5.7\times10^{-4}$.

\subsection{The Screening Rate $\sK$ and its UV Correction}

The screening rate $\sK$ is the fractional dilution of the dark-energy
\emph{pressure} per e-fold. It follows from the continuity equation and is
fixed by the equation of state alone,
\begin{equation}
  \sK \;=\; -3w_0(1+w_0) \;=\; 3\,s_{\rm DE}\,s_m \;=\; \val{S_K_d6}\,,
  \label{eq:sKdef}
\end{equation}
with $s_{\rm DE} = |w_0| = \val{S_DE_d6}$ and $s_m = 1+w_0 = \val{S_M_d6}$. It contains
no factor of $H$ and no density, which is why the screening fraction built from
it is a purely geometric lens (Paper~9). It is \emph{not} the Bellini--Sawicki
kineticity $\alpha_K$ \citep{BelliniSawicki2014,Cheung2008,Gubitosi2013}, which
for the SSEE action has the distinct value $\alpha_K = \val{alpha_K_z0_d6}$
\citep{SSEE7}. In units $\rhoc = 1$ the same number $\sK$ is reproduced by
$6\left(X K_X + 2X^2 K_{XX}\right)$ evaluated on the saturated background, and
that form is what propagates the UV correction below.
With $K_{XX} = 2/M^4$ and the self-consistent $X_{\rm bg}^{\rm UV}$:
\begin{equation}
  \sK^{\rm full}
  \;=\; \sK^{\rm IR}
        + \underbrace{\frac{24\,(X_{\rm bg}^{\rm UV})^2}{M^4\,\rhoc}}_{\rm UV\;corr.}
  \;=\; \val{S_K_d5} + 0.01361
  \;=\; \val{S_K_FULL_d5}\,.
  \label{eq:sKfullform}
\end{equation}
The UV correction is $+3.37\%$ over the IR value.

\subsection{Ghost-free and gradient-stability conditions}
\label{subsec:stability_UV}

For a $K(X)$ theory, the two fundamental stability conditions are
\citep{Garriga1999}:
\begin{enumerate}
  \item \textbf{No-ghost (positive kinetic energy):} $K_X > 0$.
  \item \textbf{Gradient stability (no gradient instability):} $c_s^2 > 0$,
    equivalently $K_X + 2XK_{XX} > 0$.
\end{enumerate}

For $K(X) = X/\KAL + X^2/M^4$ with $\KAL, M^4 > 0$ and $X \geq 0$:
\begin{equation}
  K_X = \frac{1}{\KAL} + \frac{2X}{M^4} > 0 \quad \forall\, X \geq 0.
  \label{eq:ghost_free_UV}
\end{equation}
The no-ghost condition holds \emph{for all} $X \geq 0$ in both IR and UV regimes.
\begin{equation}
  K_X + 2X K_{XX}
  = \frac{1}{\KAL} + \frac{2X}{M^4} + \frac{4X}{M^4}
  = \frac{1}{\KAL} + \frac{6X}{M^4}
  > 0 \quad \forall\, X \geq 0.
  \label{eq:gradient_stable_UV}
\end{equation}
The gradient stability condition also holds globally.
Together, equations~(\ref{eq:ghost_free_UV})--(\ref{eq:gradient_stable_UV}) imply
$c_s^2 = K_X/(K_X + 2XK_{XX}) \in (1/3,\,1)$ for all $X > 0$, with:
\begin{alignat}{2}
  X \to 0:    &\quad c_s^2 \to 1 \quad &\text{(IR limit)},\\
  X \to \infty: &\quad c_s^2 \to 1/3 \quad &\text{(UV conformal limit)}.
\end{alignat}
At the cosmological background $\varepsilon \equiv X_{\rm bg}^2/M^4 \approx 5.7\times10^{-4}$,
the full K(X) formula gives $c_s^2 \approx 0.967$ (see \S\ref{sec:consistency}), well above $1/3$.
The UV completion therefore inherits the full ghost-free and gradient-stable
structure of the IR theory, with no new instabilities introduced at any energy
below the UV cutoff $M^4$.

\begin{figure}[ht]
  \centering
  \includegraphics[width=0.85\textwidth]{../results/figures/fig_paper10_KX_profile.pdf}
  \caption{%
    Kinetic function $K(X)/X = 1/\mathrm{KAL}_0 + X/M^4$ vs the dimensionless
    kinetic ratio $u \equiv X/M^4$.
    At the cosmological background ($u_{\rm bg} = X_{\rm bg}/M^4 \approx 1.55\times10^{-3}$,
    dash-dotted; equivalently $\varepsilon \equiv X_{\rm bg}^2/M^4 \approx 5.7\times10^{-4}$
    in $\rhoc$ units), the IR term dominates and the Hubble flow is in the linear
    kinetic regime.
    Above $u \approx 1$ (green dashed), the $X^2/M^4$ UV term dominates
    and Vainshtein suppression activates (see Paper~8).
  }
  \label{fig:KX_profile}
\end{figure}

% ═══════════════════════════════════════════════════════════════════════════════
\section{UV Completion of the Hubble Tension Resolution}
\label{sec:hubble}
% ═══════════════════════════════════════════════════════════════════════════════

\subsection{Screening fraction and local $H_0$}

The screening fraction (Paper~9) generalises in the UV to
\begin{equation}
  \fsc^{\rm UV}
  \;=\; \frac{\sK^{\rm full}}{3\,\MIRA}
  \;=\; \frac{\val{S_K_FULL_d5}}{3\times \val{MIRA_d5}}
  \;=\; 0.06952\,,
  \label{eq:fUV}
\end{equation}
a $+3.37\%$ increase over the IR value $\fsc^{\rm IR} = 0.06725$.
The local Hubble constant is
\begin{equation}
  H_0^{\rm glob,UV}
  \;=\; H_0^{\rm SH0ES}\,(1 - \fsc^{\rm UV})
  \;=\; 73.04\times(1 - \val{F_SCREEN_d7})
  \;=\; 67.96214 \text{ km/s/Mpc}\,,
  \label{eq:H0UV}
\end{equation}
the headline UV result: it differs from the pure number
$3(\phiG+\pi)^2 = 67.96214$ by $4.2\times10^{-6}$, refining the IR residual
$0.166$ ($0.17\sigma$) by a factor of $4\times10^4$ with the same input.

\paragraph{What this agreement does and does not constrain.}
The SH0ES uncertainty propagates to $\pm0.968$ km/s/Mpc on the output, which is
far larger than the $4.2\times10^{-6}$ residual. Varying the cutoff over
$M^4 \in [0.2, \infty)\times 45\alpha^2\rhoc$ moves the output by at most
$0.57$ km/s/Mpc, so the cascade is compatible with that entire range: it
\emph{excludes} a cutoff ten times smaller, but it does not measure $M^4$.
Nor does it fix the pair $(1/\KAL, 1/M^4)$ separately --- a $1\%$ increase in
$\KAL$ is compensated by a $2.01\%$ increase in $M^4$. The value
$M^4 = 45\alpha^2\rhoc$ is therefore an input from the inflationary sector
($\alpha=\phiG^4/3$, Paper~1), tested for consistency here and not derived
from the Hubble cascade.

\subsection{Comparison with SH0ES}

\begin{table}[ht]
\centering
\caption{%
IR vs.\ UV global rate returned by the cascade from the measured local
  rate, compared with the dimensionless $3(\phiG+\pi)^2 = 67.96214$
  (Postulate~D). The UV completion with $M^4 = 45\alpha^2\rhoc$ sharpens the
  residual from $0.17\sigma$ (IR) to $4.3\times10^{-6}\sigma$ (UV),
  conditional on Postulate~C.1 (see §\ref{sec:theorem}). The propagated SH0ES
  uncertainty ($\pm0.970$ km/s/Mpc in the IR regime, $\pm0.968$ with the full
  screening fraction) dominates both.
}
\label{tab:H0}
\begin{tabular}{lllll}
\toprule
Regime & $\sK$ & $\fsc$ & $H_0^{\rm glob}$ [km/s/Mpc] & vs.\ $3(\phiG+\pi)^2$ \\
\midrule
IR ($M\to\infty$, Papers 1--9) & \val{S_K_d5} & 0.06725 & 68.13 & $0.17\sigma$ \\
UV ($M^4=45\alpha^2\rhoc$, this paper) & \val{S_K_FULL_d5} & 0.06952 & \textbf{67.96214} & $4.3{\times}10^{-6}\sigma$ (cond.$^{a}$) \\
\midrule
SH0ES \citep{Riess2022} (\emph{input}) & --- & --- & $73.04 \pm 1.04$  & --- \\
Freedman et al.\ \citep{Freedman2024} (input) & --- & --- & $69.96 \pm 1.54$ & $1.9\sigma$ (UV) \\
\bottomrule
\end{tabular}
\end{table}

Table~\ref{tab:H0} compares the IR and UV results.
Under Postulate~C.1 (§\ref{sec:theorem}), $M^4 = 45\alpha^2\rhoc$ gives, from
the SH0ES datum, $H_0^{\rm glob,UV} = \val{H0_GLOBAL_d5}$ km/s/Mpc --- the headline UV
result, refining the IR residual $0.166$ ($0.17\sigma$) to
$4.2\times10^{-6}$, and reproducing $3(\phiG+\pi)^2$ to the seventh digit.

\paragraph{Comparison with TRGB/JWST.}
The TRGB+JWST measurement of Freedman et al.\ \citep{Freedman2024} yields
$H_0 = 69.96\pm1.54$\,km/s/Mpc. Running the same lens on that datum instead of
SH0ES returns $H_0^{\rm glob} = 65.11$ (UV) and $65.26$ (IR) km/s/Mpc, i.e.\
$\sim2.0\sigma$ and $\sim1.9\sigma$ \emph{below} $3(\phiG+\pi)^2$.
SSEE therefore aligns with the SH0ES branch of the tension, not with
TRGB/JWST.  Whether this reflects a physical distinction between distance-ladder rungs
or a systematic in either measurement is beyond the scope of this framework.

\footnotesize{$^{a}$ Conditional on Postulate C.1 (UV-IR $\varphi/\pi$ separability);
not an independent derivation from first principles.}

\subsection{The UV ladder}

The chain of derivations that leads from the golden ratio to the SH0ES-consistent
value is summarised in Table~\ref{tab:ladder}.

\begin{table}[ht]
\centering
\caption{The UV ladder: $\phiG \to \alpha \to M \to \sK^{\rm full} \to H_0$.}
\label{tab:ladder}
\begin{tabular}{llll}
\toprule
Step & Expression & Value & Source \\
\midrule
Seed & $\phiG = (1+\sqrt{5})/2$ & \val{PHI_d8} & --- \\
$\alpha$-attractor & $\alpha = \phiG^4/3$ & \val{ALPHA_ATT_d8} & Paper 1 \\
UV cutoff & $M^4 = 45\alpha^2\rhoc$ & $\val{_M4_UV_d2}\,\rhoc$ & Eq.~\eqref{eq:M4main} \\
$M$ in meV & $M = \phiG^2\,5^{1/4}\,\rhoc^{1/4}$ & 9.68 meV & Eq.~\eqref{eq:Mmev} \\
$\sK^{\rm full}$ & quadratic + cutoff & \val{S_K_FULL_d5} & Eq.~\eqref{eq:sKfullform} \\
$\fsc^{\rm UV}$ & $\sK^{\rm full}/(3\MIRA)$ & 0.06952 & this paper \\
SH0ES (\emph{input}) & $73.04 \pm 1.04$ km/s/Mpc & (obs.) & Riess 2022 \\
$H_0^{\rm glob,UV}$ & $H_0^{\rm SH0ES}(1-\fsc^{\rm UV})$ & \textbf{\val{H0_GLOBAL_d5}} & Eq.~\eqref{eq:H0UV} \\
$3(\phiG+\pi)^2$ (pure, target) & --- & $\val{H0_ALG_d5}$ & Paper~1 \\
\bottomrule
\end{tabular}
\end{table}


% ═══════════════════════════════════════════════════════════════════════════════
\section{Cross-Epoch Falsification: Inflation $\leftrightarrow$ Dark Energy}
\label{sec:cross}
% ═══════════════════════════════════════════════════════════════════════════════

The central prediction of this paper is a cross-epoch consistency test.
If $\alpha = \phiG^4/3$ is the correct inflationary curvature, then
simultaneously:
\begin{enumerate}
  \item LiteBIRD (2032) should measure $r \approx 0.00813$
        [from Eq.~\eqref{eq:r}].
  \item The cascade run on the measured local rate should return
        $H_0^{\rm glob} = \val{H0_GLOBAL_d5}$ km/s/Mpc, reproducing $3(\phiG+\pi)^2$
        [Eq.~\eqref{eq:H0UV}].
\end{enumerate}

These are \emph{independent observables} at vastly different epochs
($z \sim 10^{60}$ vs.\ $z \lesssim 0.15$) linked only through $\alpha$.
A single parameter failure falsifies both predictions simultaneously.

\paragraph{Quantitative falsification criteria.}
If LiteBIRD measures $r_{\rm obs}$, the implied $\alpha_{\rm obs}
= r_{\rm obs} N^2/12$, and one predicts
\begin{equation}
  H_0^{\rm local}(\alpha_{\rm obs})
  \;=\; \frac{H_0}{1 - \sK^{\rm full}(\alpha_{\rm obs})\,/\,(3\MIRA)}\,.
  \label{eq:Hpred}
\end{equation}
If this differs from SH0ES by more than $1\sigma$ (i.e., if
$|H_0^{\rm local}(\alpha_{\rm obs}) - 73.04| > 1.04$ km/s/Mpc),
the UV-completion proposal is falsified at $1\sigma$.

\begin{figure}[ht]
  \centering
  \includegraphics[width=0.85\textwidth]{../results/figures/fig_paper10_alphaK_vs_alpha.pdf}
  \caption{%
    \textit{Top:} Full screening rate $s_K^{\rm full}(\alpha)$
    as the inflationary curvature $\alpha$ varies.
    At $\alpha = \varphi^4/3 = \val{ALPHA_ATT_d6}$ (green dashed), $s_K^{\rm full}=\val{S_K_FULL_d5}$
    (green dot); the IR limit $\val{S_K_d5}$ is shown dash-dotted blue.
    \textit{Bottom:} Global background rate $H_0^{\rm glob}(\alpha)$ returned by
    the cascade Eq.~(\ref{eq:Hpred}) from the measured
    $H_0^{\rm SH0ES}=73.04\pm1.04$ km/s/Mpc (solid green); the dotted blue line
    is the dimensionless target $3(\varphi+\pi)^2=\val{H0_ALG_d5}$ and the shaded band
    is the propagated SH0ES uncertainty $\pm0.968$ km/s/Mpc.
    If LiteBIRD measures $r_{\rm obs}$, one reads off $\alpha_{\rm obs}=r_{\rm obs}N^2/12$
    and the $x$-axis gives the falsifiable output.
    At $\alpha=\varphi^4/3$ the cascade returns $67.96214$ km/s/Mpc, i.e.\
    $4.2\times10^{-6}$ from the target. Note that the band is far wider than
    that residual: the agreement is not a measurement of $\alpha$.
  }
  \label{fig:cross_epoch}
\end{figure}

% ═══════════════════════════════════════════════════════════════════════════════
\section{Towards a First-Principles Derivation}
\label{sec:firstprinciples}
% ═══════════════════════════════════════════════════════════════════════════════

The identity $M^4 = 45\alpha^2\rhoc$ has been \emph{found} numerically and
verified algebraically (Section~\ref{sec:alpha}).
Its \emph{derivation} from first principles---without using SH0ES as
input---remains under investigation.
We outline three candidate routes.

\subsection{Route~A: Lagrangian matching with $K_\alpha(X)$}

The $\alpha$-attractor kinetic Lagrangian is
\begin{equation}
  K_\alpha(X) \;=\; -3\alpha\ln\!\left(1 - \frac{X}{3\alpha}\right)
  \;\approx\; X + \frac{X^2}{6\alpha} + \frac{X^3}{27\alpha^2} + \cdots
\end{equation}
If the SSEE Lagrangian $K(X) = X/\KAL + X^2/M^4$ represents the first two
terms of $K_\alpha(X)/\KAL$ (after normalising the kinetic term to
$X/\KAL$), the matching condition is
\begin{equation}
  \frac{1}{M^4} \;=\; \frac{1}{6\alpha\,\KAL^2}\,,
  \qquad\Rightarrow\qquad
  M^4_{\rm A} \;=\; 6\alpha\,\KAL^2 \;=\; 2\phiG^4\,\KAL^2
  \;\approx\; 418\,\rhoc\,.
  \label{eq:MA}
\end{equation}
Route~A gives $M^4_{\rm A} \approx 418\,\rhoc$, a factor $\approx 1.78$ above
the target $5\phiG^8\rhoc \approx 234.9\,\rhoc$.
The discrepancy is
\begin{equation}
  \frac{M^4_{\rm A}}{M^4_{\rm UV}} \;=\;
  \frac{2\phiG^4\,\KAL^2}{5\phiG^8}
  \;=\; \frac{(\phiG+3\pi)^2}{10\phiG^4}
  \;\approx\; 1.779\,,
\end{equation}
which does not simplify to a known SSEE constant.
Route~A thus requires an additional constraint to close---either a
field-redefinition factor or an additional normalisation condition inherent
to the $\alpha$-attractor sector.

\subsection{Route~B: Technical naturalness}

An EFT argument based on technical naturalness \citep{tHooft1980} demands
that the UV operator $X^2/M^4$ does not destabilise the IR physics.
Requiring that the UV correction to $\sK$ is of order $\sK^{\rm IR}$
itself gives a naturalness scale
\begin{equation}
  \delta\sK \;\sim\; \sK^{\rm IR}
  \;\Rightarrow\;
  \frac{24\,X_{\rm bg}^2}{M^4} \;\sim\; \sK^{\rm IR}
  \;\Rightarrow\;
  M^4 \;\sim\; \frac{24\,X_{\rm bg}^2}{\sK^{\rm IR}}\,.
\end{equation}
Evaluating with the IR background: $X_{\rm bg}^{\rm IR} = \sK^{\rm IR}\KAL/6$,
\begin{equation}
  M^4_{\rm B}
  \;\sim\; \frac{24\,(\sK^{\rm IR}\KAL/6)^2}{\sK^{\rm IR}}
  \;=\; \frac{2\,\sK^{\rm IR}\,\KAL^2}{3}
  \;=\; M^4_{\rm r2} \;\approx\; 4.1\,\rhoc\,.
\end{equation}
Route~B gives the self-coupling scale $M^4_{\rm r2}$, not $5\phiG^8\rhoc$.
The correct $M^4 = 5\phiG^8\rhoc$ corresponds to the regime where the UV
correction is $\approx 3.37\%$ of $\sK^{\rm IR}$---a \emph{sub-unit}
naturalness condition.

\subsection{Route~C: Full $P(X,\phi)$ Lagrangian and Conjecture~C.1}
\label{sec:routeC}
\label{sec:theorem}

\paragraph{Formal conjecture.}
The algebraic identities of Section~\ref{sec:alpha} motivate:

\medskip
\begin{quote}
\textbf{Conjecture~C.1 (Quintessential UV scale).}
\textit{In the SSEE framework, the UV coefficient of the screening
sector satisfies
\begin{equation}
  \frac{\partial^2 K}{\partial X^2}\bigg|_{X=0}
  \;=\; \frac{2}{M^4}
  \;=\; \frac{2}{45\alpha^2\,\rhoc}
  \;=\; \frac{2}{5\phiG^8\,\rhoc}\,,
  \label{eq:conjecture}
\end{equation}
where $\alpha = \phiG^4/3$ is the $\alpha$-attractor curvature fixed by
the inflationary sector and $\rhoc = 3\Mpl^2 H_0^2$ uses the algebraic anchor
$H_0^{\rm alg}/(\mathrm{km\,s^{-1}\,Mpc^{-1}}) = 3(\phiG+\pi)^2$ (Paper~1).
No low-redshift distance-ladder input is used. The statement is about the
cutoff $M$ of Eq.~\eqref{eq:KX}, not about the coefficients of the
Paper~7 condensate; relating the two functionals is a separate open
question.}
\end{quote}

\medskip
Conjecture~C.1 is promoted below to a conditional theorem (Conditional
Theorem~C.1), under the UV-IR separability postulate; the UV-completion chain
then closes within the SSEE algebraic structure:
\begin{equation}
  \phiG \;\xrightarrow{\alpha = \phiG^4/3}\; \alpha
        \;\xrightarrow{M^4 = 45\alpha^2\rhoc}\; M
        \;\xrightarrow{\sK^{\rm full}}\; \fsc
        \;\xrightarrow{\times H_0^{\rm SH0ES}}\; H_0^{\rm glob}
  \;=\; 67.96214\;\text{km/s/Mpc}\,.
\end{equation}
Conditional on that postulate, the cascade reproduces $3(\phiG+\pi)^2$ to
$4.2\times10^{-6}$. The agreement is a coincidence of central values and not a
measurement of the cutoff: the propagated SH0ES uncertainty is $\pm0.968$
km/s/Mpc, which covers $M^4$ from $0.2\times$ its nominal value upwards. The
result is not unconditional.

\paragraph{Diagnosing the Route~A discrepancy.}
Route~A matches $K(X)/\KAL$ to the Taylor expansion of $K_\alpha(X)/\KAL$
and obtains $M^4_{\rm A} = 6\alpha\,\KAL^2 = 2\phiG^4\,\KAL^2 \approx 418\,\rhoc$.
The discrepancy with the UV value is
\begin{equation}
  \frac{M^4_{\rm A}}{M^4_{\rm UV}}
  \;=\; \frac{2\phiG^4\,\KAL^2}{5\phiG^8}
  \;=\; \frac{2\KAL^2}{5\phiG^4}
  \;=\; \frac{(\phiG+3\pi)^2}{10\phiG^4}
  \;\approx\; 1.779\,.
  \label{eq:routeA_ratio}
\end{equation}
The factor $(\phiG+3\pi)^2/(10\phiG^4)$ does not reduce to a known SSEE
constant, which indicates that Route~A uses the wrong normalisation for
the kinetic-term matching.

The physical resolution is expected in the \emph{field redefinition} at the
inflaton-to-quintessence transition.
In quintessential inflation, the canonical inflaton $\chi$ is related to $\phi$
via $\phi = f(\chi)$, where $f$ carries $\alpha$-attractor factors
(e.g., $f(\chi) = \sqrt{6\alpha}\tanh(\chi/\sqrt{6\alpha})$ for the
Starobinsky--Higgs embedding).
When expressed in terms of $\phi$ (the late-time dark energy field), the
second-order kinetic coefficient picks up a factor from $|df/d\chi|^2$
evaluated at the transition, which would shift $M^4_{\rm A}$ by
$(10\phiG^4)/(\phiG+3\pi)^2 \approx 0.562$ --- precisely the inverse of
the Route~A discrepancy.

\paragraph{UV-IR separability and the conditional derivation.}

We now show that Conjecture~C.1 follows from one additional postulate,
which converts it into a conditional theorem.

The $\alpha$-attractor kinetic Lagrangian is
\begin{equation}
  K_\alpha(X) \;=\; -3\alpha\ln\!\!\left(1 - \frac{X}{3\alpha}\right)
  \;=\; X \;+\; \frac{X^2}{6\alpha} \;+\; \frac{X^3}{27\alpha^2} \;+\; \cdots
  \label{eq:Kalpha_taylor}
\end{equation}
with $\alpha = \phiG^4/3$.
In particular, the second derivative at the origin,
\begin{equation}
  \frac{\partial^2 K_\alpha}{\partial X^2}\bigg|_{X=0} \;=\; \frac{1}{3\alpha}\,,
  \label{eq:Kalpha_XX}
\end{equation}
depends \emph{only} on $\alpha$ --- and hence only on $\phiG$.

The SSEE dark-energy basis is defined by the first-order kinetic coefficient
$K_X^{\rm DE}|_0 = 1/\KAL$, where $\KAL = (\phiG+3\pi)/2$ encodes the
late-time equation of state.
A field-argument rescaling
\begin{equation}
  K_{\rm DE}(X) \;=\; K_\alpha\!\left(\frac{X}{\KALeff}\right)
  \label{eq:KDE_rescaling}
\end{equation}
reproduces $K_X^{\rm DE}|_0 = 1/\KALeff$ and yields a second-order
coefficient
\begin{equation}
  \frac{\partial^2 K_{\rm DE}}{\partial X^2}\bigg|_{X=0}
  \;=\; \frac{1}{\KALeff^2}\,\frac{\partial^2 K_\alpha}{\partial X^2}\bigg|_{X=0}
  \;=\; \frac{1}{3\alpha\,\KALeff^2}
  \;=\; \frac{2}{M^4}\,,
  \label{eq:M4_KALeff_deriv}
\end{equation}
so
\begin{equation}
  \boxed{M^4 \;=\; 6\alpha\,\KALeff^2\,.}
  \label{eq:M4_KALeff}
\end{equation}
Route~A is the special case $\KALeff = \KAL$, giving
$M^4_{\rm A} = 6\alpha\,\KAL^2 \approx 418\,\rhoc$
--- too large by the factor $(\phiG+3\pi)^2/(10\phiG^4)\approx 1.779$
derived in Eq.~\eqref{eq:routeA_ratio}.

\medskip
\noindent\textbf{No-$\pi$ separability postulate.}
The UV cutoff $M$ is set at the inflationary scale, before the dark-energy
vacuum has established itself.
At that scale, the only algebraic information available is the inflationary
$\alpha$-attractor parameter $\alpha = \phiG^4/3$.
The dark-energy constant $\KAL = (\phiG+3\pi)/2$ --- which carries $\pi$
through the late-time equation of state $w_0 = -\mathrm{AURA}/\Omega$
(the CPL parameterization \citep{ChevallierPolarski2001,Linder2003}) ---
is an \emph{infra-red} quantity absent at the transition.

We therefore postulate:
\begin{quote}
\textit{(Separability)} The effective normalisation $\KALeff$ in
Eq.~\eqref{eq:M4_KALeff} is a polynomial in $\phiG$ alone; it does not
involve $\pi$.
\end{quote}

Under this postulate, $\KALeff^2$ must be a polynomial in $\phiG$ with rational
coefficients.
The constraint Eq.~\eqref{eq:M4_KALeff} then reads
\begin{equation}
  \KALeff^2 \;=\; \frac{M^4}{6\alpha} \;=\; \frac{M^4}{2\phiG^4}\,.
  \label{eq:KALeff_constraint}
\end{equation}
The unique monomial in $\phiG$ satisfying both Eq.~\eqref{eq:KALeff_constraint}
and the empirically established $M^4 = 5\phiG^8\,\rhoc$ is
\begin{equation}
  \KALeff^2 \;=\; \frac{5\phiG^8}{2\phiG^4} \;=\; \frac{5}{2}\,\phiG^4
              \;=\; \frac{15}{2}\,\alpha\,,
  \qquad
  \KALeff \;=\; \phiG^2\sqrt{\tfrac{5}{2}}\,.
  \label{eq:KALeff_result}
\end{equation}
Note that $\KALeff^2 = (5/2)\phiG^4\,\rhoc$ is indeed a power of $\phiG$ times
a rational, with no $\pi$, confirming the postulate is self-consistent. The
factor $\rhoc$ is carried by $\KALeff^2$ and not by the rational: the
dimensionless content of the whole construction is the single number
$M^4/(2\phiG^4\rhoc) = 5/2$.
Substituting back into Eq.~\eqref{eq:M4_KALeff}:
\begin{equation}
  M^4 \;=\; 6\alpha \times \tfrac{15}{2}\,\alpha
      \;=\; 45\alpha^2\,\rhoc
      \;=\; 5\phiG^8\,\rhoc\,.
  \label{eq:M4_conditional}
\end{equation}
This is precisely Conjecture~C.1.

The Route~A discrepancy is now transparent. Since $M^4 \propto \KALeff^2$, a
ratio $r$ in the normalisation appears as $r^2$ in the cutoff:
$\KAL/\KALeff = \val{KAL_KALeff_d6}$ and $(\KAL/\KALeff)^2 = \val{KAL_KALeff2_d6}$, which is exactly
the factor by which Route~A overshoots ($\val{seis_alfa_KAL2_d2}\,\rhoc$ against
$\val{_M4_UV_d2}\,\rhoc$). The discrepancy is thus $33\%$ in the normalisation and
$78\%$ in the cutoff: the same discrepancy seen at two powers. This ratio is
the field-redefinition factor --- the ``missing'' $\pi$-dependent normalisation
between the inflationary and dark-energy kinematic bases. It is taken in units
$\rhoc = 1$; $\KAL$ is dimensionless while $\KALeff$ carries $\rhoc^{1/2}$, so
the two are not like-dimensioned objects and $\val{KAL_KALeff_d6}$ must not be read as a
comparison between them.

\medskip
\begin{quote}
\textbf{Conditional Theorem~C.1.}
\textit{If the effective kinetic normalisation $\KALeff$ at the
inflaton-to-quintessence transition is polynomial in $\phiG$ (the
UV-IR separability postulate), then the unique solution to
$M^4 = 6\alpha\,\KALeff^2$ consistent with the inflationary
$\alpha$-attractor $K_\alpha$ is
\begin{equation}
  M^4 \;=\; 45\alpha^2\,\rhoc \;=\; 5\phiG^8\,\rhoc\,.
  \label{eq:CT1}
\end{equation}
The UV ladder closes:
$\phiG \to \alpha = \phiG^4/3 \to M^4 = 5\phiG^8\rhoc
\to \sK^{\rm full} = \val{S_K_FULL_d5}
\to \fsc^{\rm UV} \to H_0^{\rm glob} = \val{H0_GLOBAL_d5}\;\mathrm{km/s/Mpc}$
(from the SH0ES datum; $4.2\times10^{-6}$ from $3(\phiG+\pi)^2$;
conditional on Postulate C.1),
entirely within the $\{\phiG,\pi\}$ algebraic structure of SSEE.}
\end{quote}

The unconditional proof requires an explicit derivation of $\KALeff$
from the $P(X,\phi)$ Lagrangian at the inflationary transition ---
specifically, showing that the field-space Jacobian $\partial\phi/\partial\chi|_{\rm transition}$
evaluates to $\phiG^2\sqrt{5/2}$ rather than $\KAL = (\phiG+3\pi)/2$.
This is deferred to Paper~B of the SSEE series.


% ═══════════════════════════════════════════════════════════════════════════════
\section{Additional UV Predictions}
\label{sec:predictions}
% ═══════════════════════════════════════════════════════════════════════════════

Beyond $H_0$, the UV completion $M^4 = 45\alpha^2\rhoc$ modifies several
EFT observables.

\paragraph{Screening rate $\sK$.}
The UV-corrected screening rate $\sK^{\rm full} = \val{S_K_FULL_d5}$ is measurable
through the equation of state, since $\sK^{\rm IR} = -3w_0(1+w_0)$.
Current constraints from DESI DR2 BAO \citep{AbdulKarim2025} set the background
context; forthcoming Euclid and DESI~Y5 data will probe $w_0$, and hence $\sK$, at the
$\sim1\%$ level, distinguishing the UV-completed theory from the IR limit by
$3.4\,\sigma_{\rm syst}$.

\paragraph{Effective equation of state.}
The UV term $X^2/M^4$ in $K(X)$ modifies the effective pressure as
\begin{equation}
  p_\phi^{\rm UV} \;=\; \frac{X}{\KAL} + \frac{X^2}{M^4} - V(\phi)\,,
  \qquad
  \rho_\phi^{\rm UV} \;=\; \frac{X}{\KAL} + \frac{3X^2}{M^4} + V(\phi)\,,
\end{equation}
Under the background matching of Eq.~\eqref{eq:quadratic} --- which holds
$\rho_\phi$ and $\rho_\phi(1+w_0)$ fixed at their IR values --- $w_0$ is
unchanged by construction (the field value $X_{\rm bg}$ adjusts instead).
A residual shift appears only if one holds the field configuration
$(X_{\rm bg}, V)$ fixed when switching on the UV term, in which case the
first-order estimate is
\begin{equation}
  \Delta w_0^{\rm UV}
  \;=\; w_0^{\rm UV} - w_0^{\rm IR}
  \;\approx\; \frac{(1-3w_0)\,X_{\rm bg}^2/M^4}{\rho_\phi}
  \;\approx\; +2.4\times10^{-3}\,.
\end{equation}
Either way, $|\Delta w_0| \lesssim 0.002$ lies well below the projected
Euclid sensitivity ($\sigma_{w_0} \approx 0.01$): the UV term leaves the
background equation of state observationally indistinguishable from the
IR limit, and the discriminating observable remains $\sK$.

\paragraph{Sound speed.}
The adiabatic sound speed of scalar-field excitations around the homogeneous
background is
\begin{equation}
  c_{s,\rm ad}^2 \;=\; \frac{K_X}{K_X + 2X K_{XX}}
          \;=\; \frac{1/\KAL + 2X/M^4}{1/\KAL + 6X/M^4}
          \;\approx\; 0.967\,.
\end{equation}
This quantity governs how fast small adiabatic excitations of $\phi$ propagate
on top of the smooth dark-energy background.
It is \emph{not} the same as the causal-perturbation speed $c_s^2 \approx 0$
reported in Paper~5, which is the Israel--Stewart bulk-viscosity transport
speed in the hydrodynamic limit of the dark-fluid sector.
The two quantities coexist without contradiction:
$c_{s,\rm ad}^2$ applies to scalar-field modes at the background level,
while Paper~5's $c_s^2 \approx 0$ is an IS closure condition that suppresses
causal bulk modes at the perturbative level.
The 3.3\% UV suppression of $c_{s,\rm ad}^2$ relative to unity is below
current CMB constraints on dark-energy perturbations, but is in principle
detectable by next-generation surveys (Euclid, CMB-S4).


% ═══════════════════════════════════════════════════════════════════════════════
\section{Consistency with the SSEE Suite}
\label{sec:consistency}
% ═══════════════════════════════════════════════════════════════════════════════

Table~\ref{tab:consistency} summarises how the UV completion is consistent
with the nine preceding papers.

\begin{table}[ht]
\centering
\caption{Consistency of the UV completion with the SSEE suite.}
\label{tab:consistency}
\resizebox{\textwidth}{!}{%
\begin{tabular}{lll}
\toprule
Paper & Result used & UV impact \\
\midrule
Paper 1 & $\alpha = \phiG^4/3$ → $r = 0.008131$ & \textbf{Direct input:} $M^4 = 45\alpha^2\rhoc$ \\
Paper 5 & IS: $c_s^2 \approx 0$ (bulk-viscosity mode) & Compatible: distinct from $c_{s,\rm ad}^2 \approx 0.967$ (field mode) \\
Paper 7 & Ghost condensate: $w_0$, $c_s^2$ from the Lagrangian & Independent: $\sK^{\rm IR}$ comes from $w_0$, not from the action \\
Paper 8 & $\MIRA = \AURA/2$; Vainshtein $r_V \propto 1/M^2$ & $r_V$ enlarged to $\sim 10^{44}$ m at UV; double GR protection \\
Paper 9 & $H_0^{\rm glob,IR} = 68.13$ (0.17$\sigma$) & UV closes to $\val{H0_GLOBAL_d5}$ ($4.3{\times}10^{-6}\sigma$; conditional$^{a}$) \\
Papers 2--4,6 & CMB, MCMC, $\sigma_8$, $f\sigma_8$ & $M$ only modifies $\sK$; the algebraic constant $\MIRA = (3\phiG+\pi)/4$ is scale-invariant and unchanged by UV flow \\
\bottomrule
\end{tabular}%
}
\footnotesize{$^{a}$ Canonical UV result $H_0^{\rm glob,UV}=\val{H0_GLOBAL_d5}$ km/s/Mpc, obtained by running
the lens on the measured SH0ES rate, conditional on Postulate C.1 (UV-IR separability,
§\ref{sec:theorem}). The propagated uncertainty is $\pm0.968$ km/s/Mpc and dominates the
$4.2\times10^{-6}$ residual; the absolute-scale origin (Postulate~D) remains open, as in $\Lambda$CDM.}
\end{table}

The AURA cancellation of Paper~9 ($\AURA$ drops out of $\sK/(3\MIRA)$
in the IR limit) is technically broken in the UV because $\sK^{\rm full}$
acquires an AURA-dependent correction.
The canonical UV result $H_0^{\rm glob,UV} = \val{H0_GLOBAL_d5}$ km/s/Mpc,
conditional on Postulate~C.1, indicates that the UV correction is a
consistent completion rather than a destabilisation of the IR identity. As stressed
throughout, the absolute dimensional scale (Postulate~D) is not an independent prediction
--- both are self-consistency checks.


% ═══════════════════════════════════════════════════════════════════════════════
\section{Summary and Conclusions}
% ═══════════════════════════════════════════════════════════════════════════════

We have shown that the SSEE dark-energy Lagrangian $K(X) = X/\KAL + X^2/M^4$
admits a unique algebraically consistent UV cutoff
\begin{equation}
  M^4 \;=\; 45\,\alpha^2\,\rhoc \;=\; 5\,\phiG^8\,\rhoc
      \;\approx\; 234.9\,\rhoc\,,
\end{equation}
where $\alpha = \phiG^4/3$ is the inflationary $\alpha$-attractor parameter of
Paper~1.
With this cutoff, the full screening rate is
$\sK^{\rm full} = \val{S_K_FULL_d5}$ and the global rate returned from the SH0ES datum is
$H_0^{\rm glob,UV} = \val{H0_GLOBAL_d5}$ km/s/Mpc, conditional on
Postulate~C.1, reproducing $3(\phiG+\pi)^2$ to the seventh digit.

The key conclusions are:
\begin{enumerate}
  \item The algebraic identity $45\alpha^2 = 5\phiG^8$ (exact, $|{\rm diff}|=0$)
        connects the inflationary curvature $\alpha$ to the dark-energy UV scale.
  \item The UV correction to $\sK$ is perturbative
        ($\varepsilon \approx 5.7\times10^{-4}$) and shifts $\sK$ by $+3.37\%$.
  \item The canonical IR Hubble residual of Paper~9
        ($H_0^{\rm glob,IR}=68.13$, $0.17\sigma$) is reduced to
        $H_0^{\rm glob,UV}=\val{H0_GLOBAL_d5}$ ($4.2\times10^{-6}$) by the UV
        completion, conditional on Postulate~C.1 (§\ref{sec:theorem}).
        The cascade does not thereby measure $M^4$: see
        §\ref{sec:hubble}.
  \item The same $\alpha$ predicts $r \approx 0.00813$ (LiteBIRD 2032),
        enabling a cross-epoch falsification test independent of Hubble tension
        measurements.
  \item A first-principles derivation of $M^4 = 45\alpha^2\rhoc$ remains
        under investigation (Section~\ref{sec:firstprinciples}), with Route~C
        (full $P(X,\phi)$ Lagrangian) identified as the most promising path.
\end{enumerate}

\paragraph{Acknowledgements.}
The author thanks Dr.\ Skylee Lee for the peer-review advice that shaped the
publication strategy of this series.

% ─── References ─────────────────────────────────────────────────────────────
\clearpage
\bibliographystyle{unsrtnat}
\begin{thebibliography}{99}

\bibitem{SSEE1}
M.~E.~Almeida Vallejo,
\emph{SSEE: Algebraic Framework, EFT, and $\alpha$-Attractor Inflation},
SSEE Paper~1 (2026).

\bibitem{SSEE7}
M.~E.~Almeida Vallejo,
\emph{The SSEE Dark-Energy Sector as a Ghost Condensate: Two-Term K-essence,
Algebraic Sound Speed, and Bellini-Sawicki Classification},
SSEE Paper~7 (2026).

\bibitem{SSEE8}
M.~E.~Almeida Vallejo,
\emph{Strong-Gravity Regime: Disformal Geodesics and $\MIRA$ Lensing},
SSEE Paper~8 (2026).

\bibitem{SSEE9}
M.~E.~Almeida Vallejo,
\emph{Hubble Tension via Algebraic Screening Fraction:
$\fsc = \sK/(3\MIRA) = (\pi-\phiG)/(\phiG+\pi)^2$},
SSEE Paper~9 (2026).

\bibitem{Riess2022}
A.~G.~Riess et al.,
\emph{A Comprehensive Measurement of the Local Value of the Hubble Constant},
ApJL \textbf{934}, L7 (2022) [arXiv:2112.04510].

\bibitem{Freedman2024}
W.~L.~Freedman et al.,
\emph{Status Report on the Chicago-Carnegie Hubble Program (CCHP): Three Independent Astrophysical Determinations of the Hubble Constant Using the James Webb Space Telescope},
ApJ \textbf{985}, 203 (2025) [arXiv:2408.06153].
[$H_0 = 69.96\pm1.05\,(\mathrm{stat})\pm1.12\,(\mathrm{sys})$\,km/s/Mpc].

\bibitem{BelliniSawicki2014}
E.~Bellini \& I.~Sawicki,
\emph{Maximal freedom at minimum cost: linear large-scale structure in general modifications of gravity},
JCAP \textbf{07} (2014) 050 [arXiv:1404.3713].

\bibitem{LiteBIRD2020}
LiteBIRD Collaboration,
\emph{LiteBIRD: a Satellite for the Studies of B-Mode Polarization and
Inflation from Cosmic Background Radiation Detection},
JLTP \textbf{200} (2020) 422 [arXiv:2202.02773].

\bibitem{Starobinsky1980}
A.~A.~Starobinsky,
\emph{A new type of isotropic cosmological models without singularity},
Phys.\ Lett.\ B \textbf{91} (1980) 99.

\bibitem{Kallosh2013}
R.~Kallosh \& A.~Linde,
\emph{Universality Class in Conformal Inflation},
JCAP \textbf{07} (2013) 002 [arXiv:1306.5220].

\bibitem{Galante2015}
M.~Galante, R.~Kallosh, A.~Linde \& D.~Roest,
\emph{Unity of Cosmological Inflation Attractors},
Phys.\ Rev.\ Lett.\ \textbf{114} (2015) 141302 [arXiv:1412.3797].

\bibitem{LiddleLeach2003}
A.~R.~Liddle \& S.~M.~Leach,
\emph{How long before the end of inflation were observable perturbations produced?},
Phys.\ Rev.\ D \textbf{68} (2003) 103503 [arXiv:astro-ph/0305263].

\bibitem{Akrami2020}
Planck Collaboration (Y.~Akrami et al.),
\emph{Planck 2018 results. X. Constraints on inflation},
A\&A \textbf{641}, A10 (2020) [arXiv:1807.06211].

\bibitem{CMBS4_2016}
K.~N.~Abazajian et al.\ (CMB-S4 Collaboration),
\emph{CMB-S4 Science Book, First Edition},
arXiv:1610.02743 (2016).

\bibitem{Garriga1999}
J.~Garriga \& V.~F.~Mukhanov,
\emph{Perturbations in k-inflation},
Phys.\ Lett.\ B \textbf{458} (1999) 219 [arXiv:hep-th/9904176].

\bibitem{tHooft1980}
G.~'t~Hooft,
\emph{Naturalness, chiral symmetry, and spontaneous chiral symmetry breaking},
in \emph{Recent Developments in Gauge Theories}, NATO ASI Series B \textbf{59} (1980) 135.

\bibitem{Perlmutter1999}
S.~Perlmutter et al.,
\emph{Measurements of $\Omega$ and $\Lambda$ from 42 high-redshift supernovae},
ApJ \textbf{517}, 565 (1999) [arXiv:astro-ph/9812133].

\bibitem{Riess1998}
A.~G.~Riess et al.,
\emph{Observational evidence from supernovae for an accelerating universe and a
cosmological constant},
AJ \textbf{116}, 1009 (1998) [arXiv:astro-ph/9805200].

\bibitem{Verde2019}
L.~Verde, T.~Treu \& A.~G.~Riess,
\emph{Tensions between the early and the late Universe},
Nat.\ Astron.\ \textbf{3}, 891 (2019) [arXiv:1907.10625].

\bibitem{Kamionkowski2023}
M.~Kamionkowski \& A.~G.~Riess,
\emph{The Hubble tension and early dark energy},
Annu.\ Rev.\ Nucl.\ Part.\ Sci.\ \textbf{73}, 153 (2023) [arXiv:2211.04492].

\bibitem{Planck2018}
Planck Collaboration,
\emph{Planck 2018 results VI: Cosmological parameters},
A\&A \textbf{641}, A6 (2020) [arXiv:1807.06209].

\bibitem{AbdulKarim2025}
M.~Abdul-Karim et al.\ (DESI Collaboration),
\emph{DESI DR2 BAO measurements},
arXiv:2503.14738 (2025).

\bibitem{Copeland2006}
E.~J.~Copeland, M.~Sami \& S.~Tsujikawa,
\emph{Dynamics of dark energy},
Int.\ J.\ Mod.\ Phys.\ D \textbf{15}, 1753 (2006) [arXiv:hep-th/0603057].

\bibitem{Weinberg2013review}
D.~H.~Weinberg et al.,
\emph{Observational probes of cosmic acceleration},
Phys.\ Rep.\ \textbf{530}, 87 (2013) [arXiv:1201.2434].

\bibitem{Horndeski1974}
G.~W.~Horndeski,
\emph{Second-order scalar-tensor field equations in a four-dimensional space},
Int.\ J.\ Theor.\ Phys.\ \textbf{10}, 363 (1974).

\bibitem{Gubitosi2013}
G.~Gubitosi, F.~Piazza \& F.~Vernizzi,
\emph{The effective field theory of dark energy},
JCAP \textbf{2013}, 032 (2013) [arXiv:1210.0201].

\bibitem{ArmendarizPicon2001}
C.~Armendariz-Picon, V.~Mukhanov \& P.~J.~Steinhardt,
\emph{Essentials of k-essence},
Phys.\ Rev.\ D \textbf{63}, 103510 (2001) [arXiv:astro-ph/0006373].

\bibitem{Cheung2008}
C.~Cheung, P.~Creminelli, A.~L.~Fitzpatrick, J.~Kaplan \& L.~Senatore,
\emph{The effective field theory of inflation},
JHEP \textbf{2008}, 014 (2008) [arXiv:0709.0293].

\bibitem{ChevallierPolarski2001}
M.~Chevallier \& D.~Polarski,
\emph{Accelerating universes with scaling dark matter},
Int.\ J.\ Mod.\ Phys.\ D \textbf{10}, 213 (2001) [arXiv:gr-qc/0009008].

\bibitem{Linder2003}
E.~V.~Linder,
\emph{Exploring the expansion history of the universe},
Phys.\ Rev.\ Lett.\ \textbf{90}, 091301 (2003) [arXiv:astro-ph/0208512].

\bibitem{Caldwell2002}
R.~R.~Caldwell,
\emph{A phantom menace? Cosmological consequences of a dark energy component with
super-negative equation of state},
Phys.\ Lett.\ B \textbf{545}, 23 (2002) [arXiv:astro-ph/9908168].

\bibitem{Ratra1988}
B.~Ratra \& P.~J.~E.~Peebles,
\emph{Cosmological consequences of a rolling homogeneous scalar field},
Phys.\ Rev.\ D \textbf{37}, 3406 (1988).

\end{thebibliography}

\end{document}
